% ===== tr-05 块1a：第 4 章 pp.27--31（原文页 27--31）=====
% [原文 p.27]
\section{环面退化、相交复形及其对偶}
我现在介绍由我自己与 Siebert 所发展的这一纲领的基本对象，其目的是在代数几何的框架下理解镜像对称。该纲领在 [29] 中宣布，并在系列文章 [31]、[32]、[33]、[22]、[34]、[28] 中得到进一步发展。

这一纲领的动机来自两个不同的方向。第一个方向（在很大程度上是我的动机）是前面几节中关于 SYZ 猜想极限形式的讨论。第二个方向出现在 Schröer 与 Siebert 的工作 [72]、[73] 中，这一工作使 Siebert 产生如下想法：Calabi--Yau 流形退化上的对数结构，可以使人把镜像对称看作是在退化的 Calabi--Yau 簇上进行的一种操作。Siebert 观察到，在组合层面上，镜像对称所交换的是与对数结构有关的数据以及与极化有关的数据。这一点在下一节我介绍对数结构时会解释得更清楚。Siebert 与我共同意识到，他所考虑的那些组合数据可以自然地编码在退化的对偶相交复形之中；而我们在上一节已经看到，这个复形看起来正是 SYZ 纤维化的底。于是，镜像对称所需的组合数据之间的互换，就对应于对偶相交复形上的一个离散 Legendre 变换。由此可以看清：这一进路为 SYZ 猜想提供了一种代数几何化。

要恰当地建立这套框架，必须考虑允许哪一类退化。它们当然应当是极大幂零的，但同一退化可以有许多不同的双有理模型。下面我们定义\emph{环面退化}的概念。环面退化这一类似乎相当受限，但它看起来是在镜像对称下封闭的最大一类退化：一个环面退化的镜像可以构造为一个环面退化。似乎没有其他自然的退化族具有这一性质。本节的大部分材料来自 [31] 的 §4。

粗略地说，Calabi--Yau 簇的环面退化是这样一种退化：其中心纤维是一族沿环面分层粘合起来的环面簇，而退化全空间在中心纤维中某个性态良好的集合 $Z$ 之外局部是环面的，并且该族局部由一个单项式给出。精确的技术性定义如下。

\begin{definition}{4.7.1}
设 $f:\mathcal{X}\to D$ 是相对维数为 $n$ 的紧合平坦族，其中 $D$ 是一个圆盘，$\mathcal{X}$ 是一个复解析空间（不必非奇异）。若下列条件成立，则称 $f$ 是 Calabi--Yau 簇的一个\emph{环面退化}：
\begin{enumerate}
\item 对 $t\neq 0$，$X_t$ 是不可约的、正规的 Calabi--Yau 簇，且只带典范奇点。（读者不妨假设 $t\neq 0$ 时 $X_t$ 光滑。）
% [原文 p.28]
% [Tutor Note: 原文写法是 \widetilde{\mathcal{X}}_0，波浪号只覆盖花体 \mathcal{X}，即它表示中心纤维 \mathcal{X}_0 的正规化，而非全空间 \mathcal{X} 的正规化；抄写时不可写成 \widetilde{\mathcal{X}_0}（后者语义不同）。]
\item 设 $\nu:\widetilde{\mathcal{X}}_0\to\mathcal{X}_0$ 是正规化，则 $\widetilde{\mathcal{X}}_0$ 是一族环面簇的不交并，导子轨迹 $C\subseteq\widetilde{\mathcal{X}}_0$ 是约化的，且映射 $C\to\nu(C)$ 非分歧、一般地为二对一。（导子轨迹是在 $\nu$ 不是同构的集合上自然定义的概形结构。）图表
\begin{center}
\begin{tikzcd}
C \arrow[r] \arrow[d] & \widetilde{\mathcal{X}}_0 \arrow[d,"\nu"] \\
\nu(C) \arrow[r] & \mathcal{X}_0
\end{tikzcd}
\end{center}
既是笛卡尔的又是余笛卡尔的。
\item $\mathcal{X}_0$ 是约化的 Gorenstein 空间，且导子轨迹 $C$ 限制到 $\widetilde{\mathcal{X}}_0$ 的每个不可约分支上，都是该分支上所有环面 Weil 除子的并。
\item 存在相对余维数 $\ge 2$ 的闭子集 $Z\subseteq\mathcal{X}$，使得 $Z$ 具有以下性质：$Z$ 不包含 $\widetilde{\mathcal{X}}_0$ 的任何环面分层在 $\nu$ 下的像；并且对任意点 $x\in\mathcal{X}\setminus Z$，存在 $x$ 的一个邻域 $\widetilde{U}_x$（在解析拓扑中）、一个 $n+1$ 维仿射环面簇 $Y_x$、$Y_x$ 上一个由单项式给出的正则函数 $f_x$，以及交换图表
\begin{center}
\begin{tikzcd}
\widetilde{U}_x \arrow[r,"\psi_x"] \arrow[d,"f|_{\widetilde{U}_x}"] & Y_x \arrow[d,"f_x"] \\
D' \arrow[r,"\varphi_x"] & \mathbb{C}
\end{tikzcd}
\end{center}
其中 $\psi_x$ 与 $\varphi_x$ 是开嵌入，且 $D'\subseteq D$。此外，$f_x$ 在 $Y_x$ 的每个环面除子上恰好消失一次。
\end{enumerate}
\end{definition}

\begin{example}{4.7.2}
取 $\mathcal{X}$ 为 $\mathbb{P}^3\times D$ 中由方程 $tf_4+z_0z_1z_2z_3=0$ 定义的子簇，其中 $D$ 是以 $t$ 为坐标的圆盘，$f_4$ 是 $\mathbb{P}^3$ 上的一般齐次四次多项式。容易看出，$\mathcal{X}$ 在轨迹
\[
\{t=f_4=0\}\cap\mathit{Sing}(\mathcal{X}_0)
\]
处奇异。由于 $\mathcal{X}_0$ 是坐标四面体，$\mathcal{X}_0$ 的奇异轨迹由 $\mathbb{P}^3$ 的六条坐标直线组成，而 $\mathcal{X}$ 在每条这样的直线上有四个奇点，共计 $24$ 个奇点。取 $Z=\mathit{Sing}(\mathcal{X})$。于是在 $Z$ 之外，投影 $\mathcal{X}\to D$ 是正规交叉的，这给出环面退化定义中的条件 (4)。容易验证其余条件也都满足。
\end{example}

给定环面退化 $f:\mathcal{X}\to D$，我们可以如下构造 $f$ 的\emph{对偶相交复形} $(B,\mathcal{P})$。这里 $B$ 是带奇点的整仿射流形，而 $\mathcal{P}$ 是 $B$ 的一个\emph{多面体分解}，即把 $B$ 分解为格多胞形。事实上，我们将把 $B$ 构造为一些格多胞形的并。具体地说，
% [原文 p.29]
% [Tutor Note: 原文 p.29 把 C(σ_x) = \{(rm,r) | m ∈ σ, r ≥ 0\} 中的 σ 印成了裸 σ；按上下文（σ_x ⊆ M_R 是格多胞形，C(σ_x) 是锥）应为 m ∈ σ_x（500 dpi 裁图 p4 p.29 x60--560/y240--305 可复核）。正文照原样排。]
设 $\mathcal{X}_0$ 的正规化 $\widetilde{\mathcal{X}}_0$ 写成环面簇 $X_i$ 的不交并 $\coprod X_i$，$\nu:\widetilde{\mathcal{X}}_0\to\mathcal{X}_0$ 为正规化映射。$\mathcal{X}_0$ 的\emph{分层}（strata）是集合
\[
\mathit{Strata}(\mathcal{X}_0)=\{\nu(S)\mid S\ \text{是某个}\ X_i\ \text{的环面分层}\}
\]
的元素。这里所谓环面分层，指的是某个 $(\mathbb{C}^*)^n$ 轨道的闭包。

设 $\{x\}\in\mathit{Strata}(\mathcal{X}_0)$ 是一个零维分层。对 $x$ 的邻域应用定义 4.7.1 的条件 (4)，可知存在环面簇 $Y_x$，使得在 $x$ 的某个邻域内，$f:\mathcal{X}\to D$ 局部同构于 $f_x:Y_x\to\mathbb{C}$，其中 $f_x$ 由一个单项式给出。而「$f_x$ 在 $Y_x$ 的每个环面除子上恰好消失一次」这一条件，正是 $Y_x$ 为 Gorenstein 的陈述；既然如此，$Y_x$ 就如同习题 4.6.1 中那样出现。事实上，取习题 4.6.1 中的 $M,N$，其中 $\operatorname{rank}M=\dim\mathcal{X}_0$。则存在格多胞形 $\sigma_x\subseteq M_{\mathbb{R}}$，使得 $C(\sigma_x)=\{(rm,r)\mid m\in\sigma,\ r\ge 0\}$ 就是定义环面簇 $Y_x$ 的锥。正如我们在习题 4.6.1 中所见，$x$ 在 $\mathcal{X}$ 中的小邻域应当向 $B$ 贡献一个 $\sigma_x$ 的拷贝，这正是我们构造的动机。现在我们可以描述如何通过粘合多胞形
\[
\{\sigma_x\mid \{x\}\in\mathit{Strata}(\mathcal{X}_0)\}
\]
来构造 $B$。

我们将在 $\mathcal{X}_0$ 的每个不可约分支本身都是正规的（从而 $\nu:X_i\to\nu(X_i)$ 是同构）的情形下做这件事。读者或许能够想象更一般的构造。在这一正规性假设下，$\sigma_x$ 的面与含有 $x$ 的 $\mathit{Strata}(\mathcal{X}_0)$ 的元素之间存在一一的包含反序对应。于是，若 $\sigma_x$ 与 $\sigma_{x'}$ 的某个面都对应于 $\mathcal{X}_0$ 的同一个分层，我们就可以把这些面等同起来。要证明这一等同可以通过一个整仿射变换来实现，需要一些论证，但这同样并不困难。

作完这些等同，就得到了 $B$。于是可以证明

\begin{lemma}{4.7.3}
若 $\mathcal{X}_0$ 是复 $n$ 维的，则 $B$ 是实 $n$ 维流形。
\end{lemma}

证明见 [31] 的命题 4.10。

至此 $B$ 还只是一个拓扑流形，由粘合格多胞形得到。令
\[
\mathcal{P}=\{\sigma\subseteq B\mid \sigma\ \text{是某个零维分层}\ x\ \text{的}\ \sigma_x\ \text{的一个面}\}.
\]
则 $\mathcal{X}_0$ 的分层与 $\mathcal{P}$ 的元素之间存在一一的包含反序对应。

剩下的只是给 $B$ 一个带奇点的仿射结构。事实上，我将描述 $B$ 上由 $\mathcal{X}_0$ 导出的更多一些结构，它特别地给出 $B$ 上一个带奇点的仿射结构。
% [原文 p.30]
% [Tutor Note: \mathcal{P} 在本节两用------既指 $B$ 的多面体分解，也指由诸零维分层对应的胞腔 $\sigma_x$ 及其面构成的集合；原文视两者为同一对象，故不区分。]
首先，对 $\tau\in\mathcal{P}$，令
\[
U_\tau:=\bigcup_{\{\sigma\in\mathcal{P}\mid \tau\subseteq\sigma\}}\operatorname{Int}(\sigma).
\]
沿 $\tau\in\mathcal{P}$ 的\emph{扇结构}是一个连续映射 $S_\tau:U_\tau\to\mathbb{R}^k$，满足
\begin{enumerate}
\item $S_\tau^{-1}(0)=\operatorname{Int}(\tau)$。
\item 若 $e:\tau\to\sigma$ 是包含关系，则 $S_\tau|_{\operatorname{Int}\sigma}$ 是到其像上的整仿射淹没。
\item 诸锥
\[
\{K_e:=\mathbb{R}_{\ge 0}S_\tau(\sigma\cap U_\tau)\mid e:\tau\to\sigma\}
\]
定义 $\mathbb{R}^k$ 中的一个有限扇 $\Sigma_\tau$。
\end{enumerate}
两个扇结构 $S_\tau,S'_\tau:U_\tau\to\mathbb{R}^k$ 若只相差 $\mathbb{R}^k$ 的一个整线性变换，则称为\emph{等价}的。

若 $S_\tau:U_\tau\to\mathbb{R}^k$ 是沿 $\tau\in\mathcal{P}$ 的扇结构，且 $\sigma\supseteq\tau$，则 $U_\sigma\subseteq U_\tau$。由 $S_\tau$ \emph{诱导}出的沿 $\sigma$ 的扇结构是复合映射
\[
U_\sigma\longrightarrow U_\tau\xrightarrow{S_\tau}\mathbb{R}^k\longrightarrow\mathbb{R}^k/L_\sigma\cong\mathbb{R}^\ell,
\]
其中 $L_\sigma\subseteq\mathbb{R}^k$ 是 $S_\tau(\sigma)$ 的线性张成。

\begin{definition}{4.7.4}
$n$ 维\emph{整热带流形}是如上的一对 $(B,\mathcal{P})$，连同对 $\mathcal{P}$ 的每个顶点 $v$ 所选定的一个扇结构 $S_v$，并具有如下性质：若 $v,w\in\tau$，则由 $S_v$ 与 $S_w$ 诱导出的沿 $\tau$ 的扇结构彼此等价。
\end{definition}

这样的数据给出 $B$ 一个带奇点的整仿射流形结构。设 $\Gamma\subseteq B$ 为 $\operatorname{Bar}(\mathcal{P})$（$\mathcal{P}$ 的第一重心重分）中那些既不含于 $\mathcal{P}$ 的极大胞腔之中、也不含 $\mathcal{P}$ 的顶点的胞腔之并。则 $B_0:=B\setminus\Gamma$ 可被
\[
\{\operatorname{Int}(\sigma)\mid \sigma\in\mathcal{P}_{\max}\}\cup\{W_v\mid v\in\mathcal{P}\ \text{是顶点}\}
\]
覆盖，其中 $W_v$ 是 $v\in B$ 的某些开邻域，含于 $U_v$ 中。我们在 $B_0$ 上定义仿射结构：令 $\operatorname{Int}(\sigma)$ 带有由 $\sigma$ 是格多胞形而给出的自然仿射结构，而 $S_v:U_v\to\mathbb{R}^n$ 限制到 $W_v$ 上是一个仿射坐标卡。

最后，关键在于：$\mathcal{X}_0$ 的结构给出 $(B,\mathcal{P})$ 上一个整热带流形结构。事实上，$\mathcal{P}$ 的每个顶点 $v$ 对应于 $\mathcal{X}_0$ 的一个不可约分支 $X_v$，而这个不可约分支是 $\mathbb{R}^n$ 中带扇 $\Sigma_v$ 的环面簇。此外，$\Sigma_v$ 的 $p$ 维锥与 $\mathcal{P}$ 中以 $v$ 为顶点的 $p$ 维胞腔之间存在一一对应，因为二者都对应于含于 $X_v$ 中的 $\mathcal{X}_0$ 的分层。于是有连续映射
\[
\psi_v:U_v\to\mathbb{R}^n
\]
% [原文 p.31]
它把 $U_v\cap\sigma$（对任意以 $v$ 为顶点、属于 $\mathcal{P}$ 的 $\sigma$）整仿射线性地映入 $\Sigma_v$ 中相应的锥。这样的映射由组合对应关系以及「在每个胞腔上整仿射线性」这一要求唯一确定。这些映射在每个顶点处定义一个扇结构。而且，这些扇结构是相容的，其含义是：若 $v,w\in\tau$，则 $U_\tau$ 上两个诱导出的扇结构等价。这是因为存在一个定义良好的扇 $\Sigma_\tau$，它定义了对应于 $\tau$ 的分层。

\begin{example}{4.7.5}
设 $f:\mathcal{X}\to D$ 是椭圆曲线退化到一个 $I_n$ 型纤维的退化。则 $B$ 是圆周 $\mathbb{R}/n\mathbb{Z}$，由 $\mathcal{P}$ 分解为 $n$ 条长度为 $1$ 的线段。
\end{example}

% [Tutor Note: 本节出现两个不同的 2×2 整矩阵，切勿混淆：例 4.7.6 的 \begin{pmatrix}1&4\\0&1\end{pmatrix} 是 $\Lambda$ 绕边重心奇点的单值变换；而 p.50 上两坐标卡之间的过渡映射是 \begin{pmatrix}1&0\\1&1\end{pmatrix}。]
\begin{example}{4.7.6}
沿用例 4.7.2，对偶相交复形是四面体的边界，其中每个面都仿射同构于标准二维单纯形，且每个顶点附近的仿射结构使得多面体分解局部看起来像 $\mathbb{P}^2$ 的扇。每条边的重心处有一个奇点，并且可以算出 $\Lambda$ 绕每个这样的奇点的单值变换在适当的基下是
\[
\begin{pmatrix} 1 & 4 \\ 0 & 1 \end{pmatrix}.
\]
\end{example}

\begin{example}{4.7.7}
考虑例 4.3.2 中的多胞形 $\Delta$。其对偶多胞形 $\nabla$ 是点 $(-1,-1,-1,-1),(1,0,0,0),\dots,(0,0,0,1)$ 的凸包。相应的射影环面簇 $\mathbb{P}_\nabla$ 有一个 crepant 消解 $X_\Sigma\to\mathbb{P}_\nabla$，其中 $\Sigma$ 是由 $\partial\Delta$ 的分解 $\mathcal{P}$ 的所有元素之上的锥组成的扇（该分解如例 4.3.2 所述）。在 $\mathbb{P}_\nabla\times\mathbb{A}^1$ 中考虑由
\[
s_0+t\sum_{m\in\nabla\cap\mathbb{Z}^4}c_ms_m=0
\]
给出的 Calabi--Yau 流形的退化族 $\mathcal{X}\to\mathbb{A}^1$，其中 $s_m$ 是 $\mathcal{O}_{\mathbb{P}_\nabla}(1)$ 中对应于 $m\in\nabla\cap\mathbb{Z}^4$ 的截面。设 $\widetilde{\mathcal{X}}$ 是 $\mathcal{X}$ 在 $X_\Sigma\times\mathbb{A}^1$ 中的真变换。则族 $\widetilde{\mathcal{X}}\to\mathbb{A}^1$ 是环面退化，其一般纤维是镜像五次超曲面，而对偶相交复形正是例 4.3.2 中构造的仿射流形 $B$。
\end{example}

对偶相交复形是不是正确的带奇点仿射流形？下面的定理为此提供了证据，并给出了这一构造与 SYZ 猜想之间的联系。

\begin{theorem}{4.7.8}
设 $\mathcal{X}\to D$ 是环面退化，对偶相交复形为 $(B,\mathcal{P})$。则存在开集 $U\subseteq B$，使得 $B\setminus U$ 收缩到 $B$ 的判别轨迹 $\Gamma$ 上；并存在 $X_t$ 的开子集 $\mathcal{U}_t$，它双全纯于 $X_\epsilon(U)$ 的一个 twist 的小形变，其中 $\epsilon=O(-1/\ln|t|)$。
\end{theorem}

这里我们不会精确说明「小形变」的含义；所谓 twist，指的是用 $B$-场扭曲 $X_\epsilon(U)$ 的复结构。更精确的陈述见 [29]；
